\section{Инженерное применение и границы}
Метод готовит правило реагирования для процессно-ориентированной информационной системы. Карточка фиксирует требования к наблюдениям, действие, порог, версию модели, диапазон калибровки, стрессовые сценарии, эффекты, стоимости и результаты сравнений. Владелец процесса утверждает действия и требования к обслуживанию; операционный контроллер регистрирует срабатывания, доступные ресурсы и последствия. Изменения маршрутов, персонала или состава задач, устойчивые ошибки прогноза либо выход эффекта за утверждённый диапазон требуют повторной калибровки. 

Допуск политики зависит от целей организации. В P0 при обеих обратных связях снижение стоимости WIP сочетается с ростом частоты затяжной перегрузки. Политики, нарушающие обязательные требования к обслуживанию, следует исключать до оптимизации стоимости. Многокритериальное предписание~\cite{multikpi} рассматривает такие компромиссы. Обоснование решения должно включать частоту тяжёлых событий и, при общих ресурсах, показатели каждого процесса наряду с общей стоимостью.

Оценка охватывает сконструированные синтетические структуры, заданные стрессы и эффекты действий. Равные веса процессов не отражают распространённость механизмов в организациях. Ресурсные соответствия и мощности задаются инженером. Реализованная процедура обнаружения предполагает полные трассы и бинарные атрибуты экземпляра для детерминированных условий; более общая маршрутизация требует более богатой модели. Корректная структура, покрытие поведения и обоснованные эффекты — отдельные условия доверия к решению. Исследования реальных журналов характеризуют предиктивный нагрузочный сигнал, а не пользу вмешательств.

Экстраполяция остаётся допущением. Калибровка P0 с нагрузочными возвратами достигает 19 незавершённых экземпляров, ниже порога реагирования 20 и критерия тяжёлого события 100. Такие журналы не различают механизмы, совпадающие в наблюдаемом диапазоне и расходящиеся при стрессе. Решения за пределами наблюдаемой нагрузки или диапазона эффектов требуют дополнительных измерений либо контролируемого пилота перед внедрением.

\section{Заключение}
Сохранение нагрузочных зависимостей улучшает подготовку реакции в последовательном контрольном случае и новой оценке при нагрузочных возвратах. С автоматическим discovery средняя нормированная стоимость снижается на 13.5\% относительно предиктивного отбора. Сохранение структуры возвратов и наблюдаемых условий маршрутизации даёт отдельный выигрыш относительно стандартного IM. Польза зависит от процесса и механизма обратной связи. Независимые исходы вмешательств наряду с предиктивной точностью и соответствием трассам дают инженерный критерий оценки моделей для подготовки реакций на перегрузку.

\begin{credits}
\subsubsection{Доступность материалов и воспроизводимость.}
Дополнительный исследовательский архив содержит протоколы, код, контрольные суммы входов, численные параметры моделей, траектории, результаты сравнений и скрипты воспроизведения анализа и таблиц. В нём различаются основной эксперимент и последующие диагностические проверки; вспомогательные исследования приведены полностью. Открытые журналы событий указаны по DOI версии 1.
\subsubsection{Использование генеративного ИИ.}
Генеративный ИИ использовался при написании кода симуляторов, анализа и графиков, организации экспериментов, проверке литературы, подготовке текста, переводе и самопроверке по артефактам. Синтетические данные получены документированными симуляторами; численные утверждения основаны на выполненных опытах и сохранённых результатах, проверенных скриптами анализа. Генеративный ИИ не входит в исполняемые модели процессов. Авторы отвечают за статью и исследовательские материалы.
\subsubsection{\discintname}
Авторы заявляют об отсутствии конфликтов интересов, относящихся к содержанию этой статьи.
\end{credits}
